\documentclass[12pt]{article}
\usepackage[utf8]{inputenc}		% Caracteres latinos
\usepackage[spanish]{babel}		% Idioma español
\usepackage{geometry}			% Organizar el documento
\usepackage{graphicx}			% Incluir gráficos
\usepackage{makecell}			% Para personalizar las celdas de una tabla
\usepackage{fancyhdr}			% Para personzalizar encabezado, pie y notas al margen
\usepackage{xcolor}				% Para darle color al texto

% Defino un color para comentarios:
\newcommand{\comment}[1]{\textcolor{gray!90}{#1}}

% Definimos la geometría de la página
\geometry{
	a4paper,                    % Tamaño del documento
	hmargin = {1.5cm, 1cm}, 	% Margen horizontal izquierdo, derecho
	vmargin = {0.5cm, 1cm},	    % Margen vertical superior, inferior
	headsep = 5mm,				% Separación entre el encabezado y el texto
	head = 3cm,					% Tamaño del encabezado
	% marginparsep = 5mm, 		% Seperación entre las notas y el texto
	% marginpar = 1.5cm,		% Tamaño de las notas
	includeall,                 % incluye el encabezado, footer y notas dentro del tamaño del documento
	nomarginpar,	            % Elimina las notas
	foot = 1.5cm,               % Tamaño del pié
	twoside,                	% Habilita el modo de impresión a doble cara
}

\selectlanguage{spanish}        % Selecciona el idioma

% Definimos el encabezado
\fancyhead[L]{
    \begin{tabular}{l l }
    	\makecell{\includegraphics[height=3cm]{logo.png}} &
    	\makecell[l]{\large\textbf{\textsc{IPEA 215} Raúl Scalabrini Ortiz} \\
    		Departamento de Matemática \\
    		Docente: Prof. Emanuel Pontoni \\
    		Alumno:\enspace\makebox[3.5in]{\hrulefill}. Curso: Primer Año ``A".}
    \end{tabular}
}

\renewcommand{\headrulewidth}{0pt}	% Establece a cero el ancho de la regla del encabezado

% Definimos el pié
\fancyfoot[C]{
	\begin{tabular}{c c}
		\makecell{\enspace\makebox[2.75in]{\hrulefill} \\ Firma del padre, madre o tutor} &
		\makecell{\enspace\makebox[2.75in]{\hrulefill} \\ Firma del alumno}
    \end{tabular}
}

\pagestyle{fancy}			% Establecemos el estilo de la página al que hemos definido con fancy.

% DOCUMENTO
\begin{document}

\begin{center}
    \textbf{PAUTAS DE TRABAJO}
\end{center}

Durante el ciclo lectivo 2019 se tendrán en cuenta los siguientes criterios:
\begin{itemize}
    \item Actitud de respeto frente al profesor y sus pares.
    \item Responsabilidad.
    \item Participación en clase.
    \item Traer las tareas completas.
    \item Cumplimiento de los materiales que se soliciten.
    \item Asistencia y puntualidad en el ingreso al aula.
    \item Cuando no se asiste a clase, debe pedir las tareas, copiar y leer para poder consultar las dudas en la clase siguiente. El alumno deberá tener una asistencia a las clases de al menos el 80\%.
    \item Los informes o trabajos prácticos deberán ser entregados en tiempo y forma, de lo contrario perderá puntaje por cada clase que se demore la entrega.
    \item Las evaluaciones se establecerán con anticipación, siendo posible la consulta sobre los temas a evaluar.
    \item Ausentarse en una evaluación implica traer certificado médico, caso contrario corresponde como calificación un 1 (uno).
    \item Cada clase es una instancia de evaluación, por lo que se llevara un registro del cumplimiento de: tareas en el hogar, materiales solicitados, trabajos áulicos, el comportamiento, la solidaridad con sus compañeros, la entrega de trabajos en tiempo y forma. El cumplimiento o no será registrado a través de signos positivos ($+$) o negativos ($-$).
    \item Está prohibido dentro del aula el uso de, teléfonos celulares, o algún dispositivo electrónico. Si así lo hiciera será retirado y dado en resguardo a dirección (salvo que el profesor lo autorice para realizar alguna actividad).
    \item Presentar el cuaderno de comunicaciones cada vez que el docente lo solicite.
\end{itemize}

%\comment{}

\medskip 
\begin{quote}
\large\textcolor{gray}{
<<Estudiar matemática, es, como enfrentarse a la ascensión de los picos más altos. En el aula se intenta proporcionar al estudiante las rutas, cuerdas e instrumental necesario para alcanzar la cima, pero tendrá que ser él quien camine y trepe. El esfuerzo no debería desanimarnos si consideramos la excepcionalidad del premio que nos aguarda al final del recorrido.>>
}
\end{quote}

\end{document}
